% Фрагмент для раздела 1 «D-разбиение»: как раскрашиваются области D_N и откуда берутся штрихи.
% Нумерация формул (4), (9) — как в основном тексте. Нужны пакеты amsmath, amssymb.
% Рисунок: Diploma/images/D-regions.png (скрипт Diploma/python/scr/main/regions.py).

\subsection*{Число корней в областях D-разбиения}

Обозначим левую часть характеристического уравнения~(4), перенесённую в одну сторону, через
\[
    P(\lambda) = \lambda + 1 - a\lambda e^{-\lambda T} - b\lambda e^{-\lambda H},
\]
а через $N(a, b)$ --- число корней $P(\lambda)$ с $\operatorname{Re}\lambda > 0$ (с учётом кратности).
Корни $P(\lambda)$ непрерывно зависят от $a$ и $b$, поэтому $N$ может измениться только в двух случаях:
\begin{enumerate}
    \item корень пересекает мнимую ось, т.\,е. $P(i\omega) = 0$ --- это и есть D-кривая~(9);
    \item корни <<приходят из бесконечности>>. Для уравнений запаздывающего типа так не бывает,
    а для уравнений нейтрального типа --- бывает, и это нужно учитывать отдельно.
\end{enumerate}

\paragraph{Разностный оператор и область $\mathbb{D}_\infty$.}
Разделим $P(\lambda) = 0$ на $\lambda$:
\[
    1 + \frac{1}{\lambda} = a e^{-\lambda T} + b e^{-\lambda H}.
\]
При $|\lambda| \to \infty$ слагаемое $1/\lambda$ исчезает, и корни с большой мнимой частью
приближаются к корням функции
\[
    D(\lambda) = 1 - a e^{-\lambda T} - b e^{-\lambda H}.
\]
При $T = 5$, $H = 1$ замена $z = e^{-\lambda}$ превращает $D$ в многочлен
\[
    D = 1 - a z^5 - b z, \qquad |z| = e^{-\operatorname{Re}\lambda}.
\]
Каждому корню $z_0$ этого многочлена соответствует целая цепочка корней $P(\lambda)$ вблизи точек
$\lambda_k = -\ln z_0 + 2\pi i k$, $k \in \mathbb{Z}$ (это следует из теоремы Руше в малых кругах вокруг $\lambda_k$),
причём $\operatorname{Re}\lambda_k \to -\ln|z_0|$ при $k \to \infty$. Поэтому:
\begin{itemize}
    \item если у многочлена есть корень с $|z_0| < 1$, то $-\ln|z_0| > 0$, и корней $P(\lambda)$ в правой
    полуплоскости \emph{бесконечно много}: $N = \infty$. Такие точки $(a, b)$ образуют область $\mathbb{D}_\infty$;
    \item если все корни лежат вне единичного круга ($|z_0| > 1$), то цепочки уходят влево,
    и $N$ конечно.
\end{itemize}
Граница $\mathbb{D}_\infty$ --- это точки, где у многочлена есть корень на окружности $|z| = 1$, т.\,е.
$D(i\omega) = 0$. Отделяя вещественную и мнимую части так же, как при выводе~(9), получаем
\emph{предельную кривую}
\[
    a(\omega) = -\frac{\sin(\omega H)}{\sin\bigl(\omega(T - H)\bigr)}, \qquad
    b(\omega) = \frac{\sin(\omega T)}{\sin\bigl(\omega(T - H)\bigr)},
\]
--- это формулы~(9) без слагаемых с $1/\omega$. Именно к ней сгущаются ветви D-кривой при $\omega \to \infty$.
Кроме того, в границу входят прямые $a + b = 1$ ($z = 1$) и $a + b = -1$ ($z = -1$), на которых
знаменатель $\sin\bigl(\omega(T - H)\bigr)$ обращается в ноль.

Точки границы $\mathbb{D}_0$, лежащие на предельной кривой, --- это и есть \emph{точки накопления}
из определения~2: при сколь угодно малом сдвиге из $\mathbb{D}_0$ в правую полуплоскость уходит
бесконечно много корней. \emph{Простые} точки границы лежат на ветвях D-кривой вне предельной кривой.

\paragraph{Направление движения корней (штриховка).}
Пусть при $(a, b)$ на D-кривой корень $\lambda = i\omega$ простой. Продифференцируем тождество
$P\bigl(\lambda(a, b), a, b\bigr) = 0$:
\[
    P'(\lambda)\,d\lambda - \lambda e^{-\lambda T} da - \lambda e^{-\lambda H} db = 0,
    \qquad
    P'(\lambda) = 1 - a e^{-\lambda T}(1 - \lambda T) - b e^{-\lambda H}(1 - \lambda H).
\]
Отсюда
\[
    \operatorname{Re} d\lambda = g_a\,da + g_b\,db, \qquad
    g_a = \operatorname{Re}\frac{i\omega e^{-i\omega T}}{P'(i\omega)}, \quad
    g_b = \operatorname{Re}\frac{i\omega e^{-i\omega H}}{P'(i\omega)}.
\]
Сдвиг в сторону вектора $(g_a, g_b)$ уводит корень $i\omega$ (а вместе с ним и сопряжённый $-i\omega$)
в правую полуплоскость. Эту сторону ветви и штрихуем: при переходе через D-кривую на заштрихованную сторону
$N$ увеличивается на~2, в обратную сторону --- уменьшается на~2. Знак $g$ постоянен вдоль ветви,
поэтому сторону достаточно определить в одной точке. Заметим, что при $a = a^* - \varepsilon$, $b = b^*$
приращение $\operatorname{Re} d\lambda = -\varepsilon g_a$ --- это вещественная часть коэффициента~(23),
умноженная на~$\varepsilon$.

Предельную кривую штрихуем в ту сторону, где у многочлена $1 - a z^5 - b z$ больше корней внутри круга
$|z| < 1$ --- туда уходит в правую полуплоскость ещё одна цепочка корней.

\paragraph{Подсчёт $N$ (принцип аргумента).}
Штриховка позволяет найти $N$ во всех областях, отсчитывая от точки $(0, 0)$, где единственный корень
$\lambda = -1$ и $N = 0$. Для проверки (и для автоматической раскраски рисунка) $N$ считается в каждой точке сетки
напрямую. Если $(a, b) \notin \mathbb{D}_\infty$, то $D(\lambda) \ne 0$ при $\operatorname{Re}\lambda \ge 0$,
и корни $P$ в правой полуплоскости совпадают с корнями функции
\[
    Q(\lambda) = \frac{P(\lambda)}{(\lambda + 1)\,D(\lambda)}
    = 1 + \frac{a e^{-\lambda T} + b e^{-\lambda H}}{(\lambda + 1)\,D(\lambda)}.
\]
При $\operatorname{Re}\lambda \ge 0$ числитель дроби по модулю не больше $|a| + |b|$, а $|D(\lambda)| \ge m$,
где $m = \min_\omega |D(i\omega)| > 0$. Поэтому $Q(\lambda) \to 1$ при $|\lambda| \to \infty$, и по принципу
аргумента (контур --- мнимая ось и большая полуокружность в правой полуплоскости, $Q(0) = 1/(1 - a - b) > 0$)
\[
    N = -\frac{1}{\pi}\,\Delta_{\omega: 0 \to +\infty} \arg Q(i\omega).
\]
При $\omega > 2(|a| + |b|)/m$ выполнено $|Q(i\omega) - 1| < 1/2$, и $Q$ уже не может обойти вокруг нуля,
поэтому приращение аргумента достаточно найти численно на конечном отрезке $\bigl[0,\, 2(|a| + |b|)/m\bigr]$.
Вблизи предельной кривой $m \to 0$, отрезок растёт, а шаг по $\omega$ приходится уменьшать --- поэтому
счёт там самый долгий.

На рис.~\ref{fig:regions} показан результат: $\mathbb{D}_0$ --- зелёная, $\mathbb{D}_2, \mathbb{D}_4, \dots$
--- от светло-жёлтого к тёмно-красному, $\mathbb{D}_\infty$ --- серая со штриховкой; в каждой области подписано $N$.
Справа увеличен верхний левый угол $\mathbb{D}_0$: здесь между простой границей $\mathbb{D}_0$ и предельной
кривой лежат узкие полосы $\mathbb{D}_2, \mathbb{D}_4, \dots$, которые сгущаются к предельной кривой.

\begin{figure}[h]
    \centering
    \includegraphics[width=\textwidth]{D-regions.png}
    \caption{D-разбиение уравнения~(4) при $T = 5$, $H = 1$. Цифрами обозначено число корней с
    $\operatorname{Re}\lambda > 0$; штрихи на границах направлены в сторону, куда корни уходят в правую полуплоскость.}
    \label{fig:regions}
\end{figure}
